\documentclass[a4paper,12pt]{article}
\usepackage{../../mypackages}
\usepackage{../../macros}

\setlength{\parindent}{0pt}


\begin{document}

\title{Chapitre 3 : Le pH}
\author{N. Bancel}
\date{Juin 2025}
\maketitle

\section{A retenir}

\section{Le pH}

\begin{compactitem}
  \item Le \textbf{pH} d’une solution renseigne sur le caractère \textbf{acide}, \textbf{basique} ou \textbf{neutre} de cette solution.
  \item C’est une grandeur sans unité, comprise entre 0 et 14, liée à la \textbf{présence des ions hydrogène} (\ce{H^+}).
  \item Si le pH est inférieur à 7, la solution est \textbf{acide} : les ions \ce{H^+} sont majoritaires.
  \item Si le pH est égal à 7, la solution est \textbf{neutre} : il y a autant d’ions \ce{H^+} que d’ions \ce{HO^-}.
  \item Si le pH est supérieur à 7, la solution est \textbf{basique} : les ions \ce{HO^-} sont majoritaires.
\end{compactitem}

\begin{figure}[H]
  \centering
  \includegraphics[width=0.8\linewidth]{img/ph_schema.jpg}
  \caption{Échelle du pH}
\end{figure}

\section*{Les transformations chimiques}

\begin{compactitem}
  \item Quand l’acide chlorhydrique et le fer sont en contact :
  \begin{compactitem}
    \item le fer et l’acide sont consommés ;
    \item il se forme du \ce{H2} (dihydrogène) et des ions fer II, \ce{Fe^{2+}}.
  \end{compactitem}

  \item Le bilan de la transformation chimique s’écrit :
  \[
  \text{acide chlorhydrique} + \text{fer} \rightarrow \ce{H2} + \text{solution de chlorure de fer (II)}
  \]

  \item Le zinc ou l’aluminium réagissent également avec l’acide chlorhydrique, contrairement au cuivre et à l’or.

  \item Les solutions acides et basiques concentrées réagissent ensemble en \textbf{dégageant de l’énergie thermique}.
\end{compactitem}

\vspace{1em}

\noindent\textbf{Remarque} : \textit{Une réaction qui libère de l’énergie thermique est qualifiée de réaction exothermique.}

\begin{figure}[H]
  \centering
  \includegraphics[width=0.5\linewidth]{img/ph_reaction.jpg}
  \caption{Échelle du pH}
\end{figure}

\section*{Sécurité}

\begin{compactitem}
  \item Les solutions acides et basiques concentrées présentent des \textbf{dangers identiques} : elles sont \textbf{corrosives}, \textbf{irritantes} et ne doivent pas être rejetées dans l’environnement.
  
  \item Pour \textbf{minimiser les risques} liés à leur utilisation, elles doivent être \textbf{diluées}, en versant toujours \textbf{la solution dans l’eau} et jamais le contraire.
\end{compactitem}

\vspace{1em}

\noindent
\textbf{Remarque} : \textit{Toute inhalation ou contact avec la peau ou les yeux peut provoquer de graves lésions. Il faut rincer immédiatement et abondamment la zone contaminée puis prévenir un médecin.}

\vspace{1em}

\noindent
\textit{Il faut manipuler dans le calme, en portant des lunettes de protection, une blouse et des gants spécifiques.}

\begin{figure}[H]
  \centering
  \includegraphics[width=0.8\linewidth]{img/secu.jpg}
\end{figure}

\end{document}
